\documentclass{article}
\usepackage{graphicx} % Required for inserting images
\usepackage{amsmath, amssymb, amsthm}
\usepackage{mathtools}
\usepackage{lettrine}

\usepackage[hidelinks]{hyperref}
\theoremstyle{definition}
\newtheorem{theorem}{Theorem}[subsection]
\newtheorem{lemma}[theorem]{Lemma}
\newtheorem{proposition}[theorem]{Proposition} 
\newtheorem{corollary}[theorem]{Corollary} 
\newtheorem{definition}[theorem]{Definition}
\newtheorem{example}[theorem]{Example}
\newtheorem{claim}[theorem]{Claim}
\newtheorem{property}[theorem]{Property}


\newcommand{\R}{\mathbb{R}}
\newcommand{\C}{\mathbb{C}}
\newcommand{\Ker}{\operatorname{Ker}}
\newcommand{\im}{\operatorname{im}}
% \newcommand{\id}{\mathrm{id}}
\newcommand{\id}{I}

 


\title{Generalized eigenvalues and Weierstrass Form}
\author{Gloria, Bodhi, Branko, Albert}
\date{April 2026}

\begin{document}

\maketitle


\tableofcontents
\newpage
\section{Invariant subspaces}
\lettrine{W}{e}
consider a linear map $T:V\rightarrow V$, where $V$ is finite dimensional. For a random subspace $S\subset V$, its image under $T$ is usually also random.
However, how about we consider those $S$ which do not ``change'' under $T$, called invariant subspaces:$$v\in S\implies T(v)\in S.$$The best way to illustrate their effectiveness is through examples\\
\subsection{Examples and properties}
\begin{example}
    The well-known subspace $$E_\lambda = \{v: T(v)=\lambda v\}$$ is invariant - each element is scaled under $T$, which is still in $E_\lambda$.\\ \\
\end{example} 


\begin{example}
    The image $$\im T=\{T(v): v\in V\}$$is invariant almost by definition.
\end{example} 

\begin{example}\label{ex:invariance1}
    The generalized kernel $$K=\{v: \exists k\in \mathbb{N}^* \text{ s.t } T^k(v)=0\}$$
is invariant under $T$. This is indeed a subspace, and the invariance here follows from picking $v\in K$ nonzero and it's corresponding $k$ such that we have 
$$T^k(v)=0\implies T^{k-1}(T(v))=0\implies T(v)\in K\ (\text{definition}).$$
\end{example}

$\newline$Since their behavior is simple enough, it is natural to search for building blocks of $V$ of invariant subspaces - similarly to how the integers are building blocks of primes. By building blocks, we mean invariant subspaces $$S_1,S_2,...,S_k$$such that for a basis $\beta_i$ for $S_i$, $$\bigcup_{i=1}^k \beta_i$$is a basis for $V$. Equivalently, $$\bigoplus_{i=1}^kS_k=V$$If we one succeeds in this, then there would be a matrix representation of $T$ in the form \[
\begin{pmatrix}
A_1 & 0   & \cdots & 0 \\
0   & A_2 & \cdots & 0 \\
\vdots & \vdots & \ddots & \vdots \\
0   & 0   & \cdots & A_k
\end{pmatrix}
\]
This structure is strictly due to the invariance of the subspaces - if $\beta_i$ is a basis for $S_i$ consisting of $\beta_{i1},\beta_{i2},...,\beta_{ik}$, then $$T(\beta_{ij})\in S_i\ \forall j\leq k, $$and therefore may be expressed only in terms of $\beta_{i1},\beta_{i2},...,\beta_{ik}$ - hence the above pattern.\\

Our goal is to achieve this decomposition with choosing the best possible invariant subspaces - so that the above $A_i$ blocks are of the simplest scheme.\\
It should already be obvious which invariant subspaces to start from. Before continuing, some properties:


\begin{property}
Invariant subspaces (except $\{\mathbf{0}\}$) $S$ under $T$ all contain an eigenvector of $T$. This is for the reason that the linear map

 $$T^*:S\rightarrow S\ \text{,} \ T^*(v)=T(v)$$   
is well defined due to the invariance. Then we just pick the eigenvector of $T^*$ (always exists).
    
\end{property}

\begin{lemma}\label{lem:invariance2}
    $$S\ \text{invariant under } T\iff S\ \text{invariant under } T-x\id\ \forall x\in \C,$$ and consequently invariant under any polynomial $P(T)$.
\end{lemma}



\section{The Jordan form}
In search of a canonical matrix representation of any $T$ -  
we begin with considering the eigenspaces. Let the eigenvalues be $\lambda_1, \lambda_2,...,\lambda_m$ for a linear operator $T$ over $V$, and the corresponding eigenspaces $$E_{\lambda_{i}}=\{v: T(v)=\lambda_i v\}$$
It is important to know that $$E_{\lambda_i}\cap E_{\lambda_j}=\{0\}  \ \text {, for all }i<j\leq k,$$and that for an arbitrary basis $\beta_i$ for $E_{\lambda_i}$, the set $$\bigcup_{i=1}^{m} \beta_i$$is linearly independent.$\newline$$\newline$
Is $\bigoplus_{i=1}^{m} E_{\lambda_i} = V$ true? - if yes, $\bigcup \beta_i$ may be our basis. 
Sadly, this is often not the case. An example is the matrix \[
A = \begin{pmatrix}
1 & 1 & 0 \\
0 & 1 & 1 \\
0 & 0 & 1
\end{pmatrix}
\]
with its corresponding linear map. If $\bigoplus E_{\lambda_i}$ is not enough to cover $V$, we want to somehow expand each $E_{\lambda_i}$ into a bigger subspace, with the importance to keep their invariance and disjointedness. $\newline \newline$

\subsection{Method}The way to this is for each $E_{\lambda_i}$ to consider $$G_{\lambda_i}=\bigcup_{i=1}^{\infty} \ker(T-\lambda_i \id)^k$$Denote $X^k_i=\ker(T-\lambda_i \id)^k$ and notice that $E_{\lambda_i}=X^1_i$. It can be seen that $$X^k_i\subseteq X^{k+1}_{i}$$ which means that for each iteration of $k$, the subspace $X^k_i$ grows. Each $X^k_i$ has 2 nice properties: it is nilpotent under $T-\lambda_i\id$ (definition), and invariant under $T-x\id\ \forall x\in \C $ which actually follows from the former. \\

\begin{lemma}
    For each $k\in \mathbb{N}, $$X^k_i$ is invariant under $T-x\id\ \forall x\in \C$
 \end{lemma}

\begin{proof}
    Follows from \ref{ex:invariance1} together with \ref{lem:invariance2}
\end{proof}

Notice that because $X^k_i\subseteq X^{k+1}_{i}\subseteq V \ \forall k$, after some iteration $N$ we must have $$X^k_i=X^{k+1}_{i}\  \forall k\geq N$$since we are bounded by the dimension of $V$.
$\\$ $\\$ Our final extendable subset is $G_{\lambda_i}=X^N_i$.

\begin{corollary}\label{col: invariance3}
    $G_{\lambda_i}$ is invariant under $T-x\id\ \forall x\in \C.$
\end{corollary}

$\newline$There is an obvious way to pick an appropriate basis for $G_{\lambda_i}$. Recall that each $v\in G_{\lambda_i}$ eventually gets mapped to $0$ under $T-\lambda_i\id$, so it is natural to look at the vectors from the sequence $$v,(T-\lambda_i\id)(v), (T-\lambda_i\id)^2(v),...$$
\begin{definition}
    A \textit{null-chain} under a map $T$ is a set of (one or more) \textbf{non-zero} vectors $$\{v_1,v_2,..v_n\}$$such that $$T(v_i)=v_{i+1}\text{ for }i\leq n-1 $$and $$T(v_n)=0$$
\end{definition}
\begin{lemma}
    A null-chain under any $T$ forms a linearly independent set.
\end{lemma}  

\begin{proof}
    Consider a null-chain $$\{v_1,v_2,...,v_n\}$$ and let 
    \begin{equation}\label{null}
        c_1v_1+c_2v_2+...+c_nv_n=0.
    \end{equation}

    
$\newline$Applying $T^{n-1}$ on both sides gives $$c_1T^{n-1}(v_1)+...+c_nT^{n-1}(v_n)=0\implies c_1v_n=0\implies c_1=0.$$
Substituting back in (\ref{null}) and applying $T^{n-2}$ on both sides gives $$c_2v_n=0\implies c_2=0.$$
Thus, inductively we obtain $$c_1=c_2=...=c_n=0$$
\end{proof} 

$\newline$So we can pick null-chain in $G_{\lambda_i}$ under $T-\lambda_i\id$, and it is of question whether it can be a basis for $G_{\lambda_i}$. It usually is not, unfortunately. For example the matrix \[
A=\begin{pmatrix}
2&1&0\\
0&2&0\\
0&0&2
\end{pmatrix}
\]has $G_2$ with dimension $3$ for the eigenvalue $2$, however the longest null-chain is of size $2$ so one is not enough.\\ \\
The trick, however, is to pick multiple such chains. In general: 
\begin{theorem}
    A vector space $V$ under a nil-potent map $T$ always has a basis composed of null-chains. 
\end{theorem} 
\textit{Proof:}
\hfill$\blacksquare$ \\ \\
Now that we can choose a nice-behaving basis $\beta_i$ for each $G_{\lambda_i}$, we proceed with our main goal: to show that these extensions of $E_{\lambda_i}$ are now enough to span the whole space. Namely:$$\bigoplus_iG_{\lambda_i}=V$$
- then $\bigcup_i\beta_i$ may be our basis for the whole space. For this to first make sense, we need to make sure that the following two theorems hold.\\ \\
\begin{theorem}\label{thrm: disj}
    $i\neq j\implies G_{\lambda_i}\cap G_{\lambda_j}=\{0\}  $
\end{theorem}
\begin{proof} The claim is equivalent to:  
\begin{align}
v\in G_{\lambda_i}\cap G_{\lambda_j}
\implies v=0 
\end{align}
which is equivalent to
\begin{align}
(T-\lambda_i I)^{k_1}(v)=(T-\lambda_j I)^{k_2}(v)=0\implies v=0.
\end{align}


where $k_1,k_2\in \mathbb{N^*}.$
Before we start, we first show the following weaker claim.

\begin{claim}\label{cl:disj}
    $i\neq j \implies E_{\lambda_i}\cap G_{\lambda_j}=\{0\}$
\end{claim}

\textit{Proof.}
    Let $u\in E_{\lambda_i}\cap G_{\lambda_j}$. Then, \begin{equation}
    \label{eq:eigen}
    T(u)=\lambda_i u
\end{equation} and $\exists k\in \mathbb{N}^*$ s.t
\begin{equation}
\label{eq:chain}
    (T-\lambda_j\id)^k(u)=0.
\end{equation}
However,
$$(\ref{eq:eigen})\implies (T-\lambda_j\id)(u)=(\lambda_i - \lambda_j)(u)\implies(T-\lambda_j\id)^k(u)=(\lambda_i-\lambda_j)^k(u) $$and therefore $$(\ref{eq:chain})\implies (\lambda_i-\lambda_j)^k(u)=0\implies u=0$$because $i\neq j$.\hfill $\blacksquare$ \\ \\
Having this, consider $u=(T-\lambda_i\id)^{k_1-1}(v)$. Then,
\begin{equation}\label{5}
(T-\lambda_i\id)^{k_1}(v)=(T-\lambda_i\id)(u)=0\implies u\in E_{\lambda_i}.
\end{equation}
However note that $$v\in G_{\lambda_j}\implies (T-\lambda_i\id)(v)\in G_{\lambda_j}$$due to \ref{col: invariance3}.
\begin{equation}\label{6}
    \implies u=(T-\lambda_i\id)^{k_1-1}(v)\in G_{\lambda_j}\implies u\in G_{\lambda_j} 
\end{equation}
Therefore (\ref{5}) and (\ref{6}) together with Claim \ref{cl:disj} imply $$u=0\implies (T-\lambda_j\id)^{k_1-1}(v)=0.$$
Similarly, letting $u_1=(T-\lambda_j\id)^{k_1-2}(v)$, the same approach implies $u_1=0$. Continuing this process inductively we arrive at $v=0.$



\end{proof}


\begin{theorem}
    For a basis $\beta_i$ of $G_{\lambda_i}$, the set $$\bigcup_i\beta_i$$is linearly independent
\end{theorem}
\begin{proof}
    From Theorem \ref{thrm: disj}, all the $G_{\lambda_i}$ are mutually disjoint (except at $0$). Therefore, it is sufficient to show that for $u_i\in G_{\lambda_i}$,
    \begin{equation}\label{8}
    u_1+u_2+...+u_m=0
    \end{equation}
    implies $u_i=0\ \forall i\leq k.$\\
    For each $u_i,\ \exists m_i\in \mathbb{N}^*$ s.t 
    \begin{equation}
        (T-\lambda_i\id)^{m_i}(u_i)=0
    \end{equation}
     Corollary \ref{col: invariance3} gives 
     \begin{equation}\label {10} 
     u_i\in G_{\lambda_i}\implies (T-x\id)(u_i)\in G_{\lambda_i}\ \forall x\in \C
     \end{equation}
     Now apply $(T-\lambda_1\id)^{m_1}$ to both sides in (\ref{8}) to vanish the first term and obtain
     \begin{equation}\label{11}
     0+(T-\lambda_1\id)^{m_1}(u_2)+...+(T-\lambda_1\id)^{m_1}(u_m)=0.
     \end{equation}
     Note that (\ref{10}) implies $(T-\lambda_1\id)^{m_1}(u_i)$ is still in $G_{\lambda_i}$ for $2\leq i\leq k$, so now we apply $(T-\lambda_2\id)^{m_2}$ to both sides in (\ref{11}) to vanish the second term, then $(T-\lambda_3\id)^{m_3}$ to vanish the third term, etc. At the end, we obtain
     \begin{equation}
         (T-\lambda_{k-1}\id)^{m_{k-1}}...(T-\lambda_2\id)^{m_2}(T-\lambda_1\id)^{m_1}(u_k)=0
     \end{equation}
     But then Theorem \ref{thrm: disj} implies $u_k=0$. Substituting back in $(\ref{8})$ and repeating the same process gives $$u_1=u_2=...=u_m=0$$
\end{proof}






























 




\end{document}
